\chapter{نتیجه‌گیری و پیشنهادها}
\label{chap:conclusion}

\section{مقدمه}
\label{sec}
در فصل‌های پیشین، ابتدا مسئله پژوهش و ضرورت استفاده از دوقلوهای دیجیتال و عامل‌های هوشمند در پایش و نگهداری پیش‌بینانه بررسی شد و سپس مبانی نظری و مطالعات مرتبط با دوقلوهای دیجیتال، سیستم‌های چندعاملی، مدل زبانی بزرگ و پیش‌بینی عمر مفید باقیمانده مورد بررسی قرار گرفت. در ادامه، معماری پیشنهادی دوقلوی دیجیتال چندعاملی طراحی و اجزای آن شامل عامل دوقلوی دیجیتال، عامل پیش‌بینی، عامل تعامل زبانی، عامل پایش و عامل \gls{Dashboard} تعریف و پیاده‌سازی شد. در فصل پنجم نیز عملکرد مدل پیش‌بینی و قابلیت‌های معماری چندعاملی از دید پیش‌بینی، تعامل، پایش و اجرای حلقه بسته تصمیم ارزیابی گردید. فصل حاضر به جمع‌بندی پژوهش و ارائه نتیجه‌گیری نهایی اختصاص دارد و در ادامه، مسیرهای پیشنهادی برای توسعه و تکمیل پژوهش در آینده ارائه می‌شوند. هدف از این فصل، قرار دادن یافته‌های فصل‌های پیشین در یک چارچوب واحد و تبیین میزان تحقق هدف پژوهش بر اساس نتایج به‌دست‌آمده است.

\section{جمع‌بندی پژوهش}
\label{sec:research-summary}
پژوهش حاضر با تمرکز بر مسئله پایش و پیش‌بینی وضعیت موتورهای توربوفن، با هدف بررسی امکان طراحی و پیاده‌سازی یک چارچوب دوقلوی دیجیتال چندعاملی شکل گرفت. مسئله اصلی از این نیاز آغاز شد که یک سامانه نگهداری پیش‌بینانه صرفاً به تولید یک مقدار عددی برای عمر مفید باقیمانده محدود نباشد، بلکه بتواند وضعیت موجودیت را به صورت یکپارچه دریافت و نمایش دهد، خروجی مدل پیش‌بینی را در اختیار اجزای دیگر قرار دهد، اطلاعات حاصل را برای کاربر قابل فهم کند و در چارچوب سیاست‌های تعریف‌شده، امکان انجام اقدامات عملیاتی را نیز فراهم آورد. بر این اساس، پژوهش به جای تمرکز صرف بر توسعه یک مدل یادگیری عمیق مستقل، بر طراحی معماری‌ای متمرکز شد که در آن مدل پیش‌بینی، دوقلوی دیجیتال و عامل‌های کاربردی از جمله زبان طبیعی، بتوانند در قالب یک سامانه واحد با یکدیگر تعامل داشته باشند.

در گام نخست، مبانی نظری مرتبط با دوقلوهای دیجیتال، سیستم‌های چندعاملی، دوقلوهای دیجیتال چندعاملی، مدل‌های زبان بزرگ و نگهداری پیش‌بینانه بررسی شد. مرور مطالعات مرتبط نشان داد که این حوزه‌ها هر یک بخش مهمی از مسئله را پوشش داده‌اند، اما ترکیب آن‌ها در قالب یک معماری اجرایی و یکپارچه همچنان دارای ظرفیت پژوهشی است. بر اساس این بررسی، معماری پیشنهادی با رویکردی ماژولار طراحی شد تا مسئولیت‌های اصلی میان عامل‌های تخصصی توزیع شوند. در این معماری، عامل دوقلوی دیجیتال مسئول نگهداری و ارائه وضعیت موجودیت، عامل پیش‌بینی مسئول تخمین \gls{RUL}، عامل پایش مسئول مشاهده و پیگیری وضعیت، و عامل تعامل زبانی مسئول برقراری ارتباط طبیعی با کاربر و هماهنگی برخی فرایندهای عملیاتی است.
در گام بعد، معماری طراحی‌شده در قالب یک نمونه آزمایشگاهی پیاده‌سازی شد. برای ارزیابی بخش پیش‌بینی، از مجموعه داده \gls{CMAPSS} و به طور مشخص زیرمجموعه \lr{FD001} استفاده شد. مدل \lr{LSTM} در این ساختار به عنوان یک مؤلفه تخصصی و قابل تعویض در نظر گرفته شد و نسخه نهایی آن با استفاده از سازوکار توجه و تابع زیان \lr{Pinball} با هدف ایجاد رفتار محافظه‌کارانه‌تر در برابر بیش‌برآورد \gls{RUL} آموزش داده شد. نتایج ارزیابی نشان داد که این مدل می‌تواند عملکرد مورد نیاز برای ایفای نقش سرویس پیش‌بینی در معماری را فراهم کند. در عین حال، جایگاه مدل در چارچوب پژوهش به گونه‌ای تعریف شد که جایگزینی آن با مدل‌های پیش‌بینی دیگر در آینده، بدون نیاز به بازطراحی اساسی سایر اجزای معماری امکان‌پذیر باشد.

در سطح سامانه، عامل زبانی طبیعی به عنوان واسط میان کاربر و عامل‌های تخصصی قرار می‌گیرد. در این ساختار، محاسبه مقدار \gls{RUL} به مدل زبانی واگذار نشده و مقدار پیش‌بینی‌شده از سرویس تخصصی دریافت می‌گردد. عامل زبانی سپس این مقدار و سایر اطلاعات موجود را برای کاربر تفسیر کرده و در مسیرهای عملیاتی تعریف‌شده، درخواست‌های مناسب را به عامل‌های مربوطه منتقل می‌کند. این جداسازی مسئولیت‌ها سبب شد که محاسبه عددی، تفسیر زبانی و کنترل عملیاتی در عامل‌ها مشخص و قابل ارزیابی باقی بمانند.

در نهایت، اجرای فرایندهای سرتاسری نشان داد که اجزای سامانه می‌توانند در یک زنجیره مشترک با یکدیگر همکاری کنند. در این زنجیره، داده وارد سامانه شده، وضعیت دوقلوی دیجیتال شکل گرفته، مدل پیش‌بینی عمر مفید را تخمین زده، عامل پایش وضعیت را ارزیابی کرده و عامل زبان طبیعی امکان تعامل کاربر با این اطلاعات و انجام فرمان‌های مجاز را فراهم کرده است. همچنین در فرایند حلقه بسته، تغییر ایجادشده در وضعیت عملیاتی دوباره به دوقلوی دیجیتال منتقل شده و در لایه پایش قابل مشاهده است. بنابراین، پژوهش حاضر نشان می‌دهد که می‌توان یک مدل پیش‌بینی و مجموعه‌ای از خدمات مستقل را در قالب یک معماری چندعاملی به یک جریان عملیاتی منسجم تبدیل کرد.

با توجه به مجموعه نتایج، می‌توان گفت مسئله پژوهش از یک دیدگاه صرفاً مدل‌محور به یک دیدگاه سامانه‌محور توسعه یافته است. در این رویکرد، کیفیت سامانه تنها به دقت مدل پیش‌بینی وابسته نیست، بلکه به توانایی اجزای مختلف در همکاری، انتقال اطلاعات، تفسیر نتایج و تبدیل خروجی تحلیلی به اقدام نیز وابسته است. از این رو، معماری پیشنهادی تلاشی برای ایجاد ارتباط میان مدل پیش‌بینی، دوقلوی دیجیتال و لایه تعامل هوشمند است و نمونه پیاده‌سازی‌شده امکان بررسی این ارتباط را در یک فرایند مشخص فراهم کرده است.



% \section{دستاوردها و نوآوری‌های پژوهش}
% \label{sec:contributions}
% مهم‌ترین دستاورد پژوهش حاضر، ارائه و پیاده‌سازی یک معماری دوقلوی دیجیتال چندعاملی است که در آن قابلیت‌های دوقلوی دیجیتال، پیش‌بینی عمر مفید باقیمانده، پایش وضعیت، تعامل زبانی و کنترل عملیاتی در قالب مجموعه‌ای از عامل‌های تخصصی به یکدیگر متصل شده‌اند. این معماری تلاش می‌کند فاصله میان یک مدل پیش‌بینی مستقل و یک سامانه تصمیم‌یار تعاملی را کاهش دهد. در ساختار پیشنهادی، هر عامل مسئولیت مشخصی دارد و ارتباط میان عامل‌ها از طریق رابط‌ها و قراردادهای تعریف‌شده انجام می‌شود. در نتیجه، معماری از یک ساختار یکپارچه و وابسته به یک مدل خاص فاصله گرفته و به سمت ساختاری ماژولار و قابل توسعه حرکت می‌کند.
% یکی از دستاوردهای اصلی پژوهش، تفکیک مسئولیت میان مدل پیش‌بینی و سایر اجزای سامانه است. مدل \lr{LSTM} در این پژوهش به عنوان هسته غیرقابل تغییر سامانه در نظر گرفته نشده، بلکه به عنوان یک مؤلفه تخصصی برای تخمین \gls{RUL} در عامل پیش‌بینی قرار گرفته است. این نحوه سازمان‌دهی باعث می‌شود مدل پیش‌بینی بتواند در آینده با مدل‌های دیگری مانند مدل‌های مبتنی بر \gls{Transformer} یا روش‌های احتمالاتی جایگزین شود، بدون آنکه منطق اصلی عامل دوقلوی دیجیتال، عامل پایش یا عامل تعامل زبانی نیازمند بازطراحی اساسی باشد. از این منظر، سهم پژوهش تنها انتخاب یک مدل یادگیری عمیق نیست، بلکه نحوه قرار دادن این مدل در یک معماری قابل توسعه است.
% دستاورد دیگر، وارد کردن ملاحظه هزینه نامتقارن خطا به زنجیره پیش‌بینی و تصمیم است. در این پژوهش، استفاده از تابع زیان \lr{Pinball} با $\tau=0.125$ با هدف ایجاد گرایش کنترل‌شده به سمت کم‌برآورد \gls{RUL} انجام شد تا احتمال پیش‌بینی بیش از حد خوش‌بینانه کاهش یابد. نتایج آزمایش نشان داد که نسبت کم‌برآورد مدل نهایی نسبت به مدل پایه افزایش یافته است. اهمیت این رویکرد در آن است که عملکرد مدل تنها بر اساس خطای عددی عمومی ارزیابی نشده و نیاز تصمیم‌گیری نگهداری نیز در طراحی آن مورد توجه قرار گرفته است. در عین حال، پژوهش نشان می‌دهد که کم‌برآورد بیش از حد نیز مطلوب نیست و انتخاب پارامتر تابع زیان باید متناسب با هزینه، ریسک و ویژگی‌های تجهیز انجام شود.
% از دیگر دستاوردهای پژوهش، استفاده از مدل زبان بزرگ در نقش یک لایه تعامل و هماهنگی، به جای استفاده از آن به عنوان جایگزین مدل تخصصی پیش‌بینی است. در معماری پیشنهادی، عامل زبانی مقدار \gls{RUL} را به صورت مستقل محاسبه نمی‌کند، بلکه اطلاعات حاصل از سرویس‌های تخصصی را دریافت و به زبان طبیعی برای کاربر تفسیر می‌کند. این تفکیک باعث می‌شود استنتاج عددی و تعامل زبانی از یکدیگر مستقل باقی بمانند و در عین حال، کاربر بتواند بدون تعامل مستقیم با رابط‌های فنی، درباره وضعیت تجهیز پرس‌وجو کند و در چارچوب فرمان‌های تعریف‌شده درخواست عملیاتی ارائه دهد.
% همچنین، پیاده‌سازی حلقه بسته پیش‌بینی، تصمیم، کنترل و بازخورد را می‌توان یکی دیگر از دستاوردهای معماری دانست. در نمونه اجرایی، خروجی پیش‌بینی صرفاً در قالب یک گزارش نمایش داده نشده، بلکه می‌تواند در مسیر تصمیم‌گیری قرار گیرد و نتیجه تصمیم پس از اعمال فرمان به وضعیت دوقلوی دیجیتال بازگردد. در نتیجه، دوقلوی دیجیتال، عامل پیش‌بینی، عامل زبان و عامل پایش در یک فرایند سرتاسری به یکدیگر متصل می‌شوند. این ویژگی، معماری را از یک سامانه صرفاً تحلیلی به سمت یک چارچوب تصمیم‌یار و عملیاتی هدایت می‌کند.
% دستاورد دیگر، ارائه یک نمونه اجرایی برای بررسی هم‌زمان چند قابلیت در یک سناریوی مشخص است. استفاده از \gls{CMAPSS} امکان ارزیابی مدل پیش‌بینی را فراهم کرده و سرویس‌های پیاده‌سازی‌شده نیز امکان بررسی تعامل میان پیش‌بینی، دوقلوی دیجیتال، پایش و عامل زبان را فراهم ساخته‌اند. از این رو، پژوهش حاضر علاوه بر تعریف معماری در سطح مفهومی، امکان بررسی رفتار آن را در سطح پیاده‌سازی و اجرای فرایندهای سرتاسری نیز فراهم کرده است.
% در مجموع، نوآوری اصلی پژوهش را می‌توان در سطح معماری و یکپارچه‌سازی قابلیت‌ها تعریف کرد. ارزش چارچوب پیشنهادی نه در ادعای ارائه یک مدل پیش‌بینی برتر نسبت به تمام روش‌های موجود، بلکه در سازمان‌دهی چند مؤلفه تخصصی در قالب یک ساختار واحد و قابل توسعه است؛ ساختاری که در آن دوقلوی دیجیتال نقش بازنمایی و نگهداری وضعیت، مدل یادگیری نقش پیش‌بینی، عامل پایش نقش مشاهده، و عامل زبان نقش تعامل و هماهنگی را بر عهده دارند. این تفکیک، زمینه را برای توسعه و جایگزینی مستقل اجزای سامانه و انتقال معماری به سناریوهای دیگر فراهم می‌کند.

\section{نتیجه‌گیری}
\label{sec}
بر اساس نتایج به‌دست‌آمده، می‌توان نتیجه گرفت که هدف اصلی پژوهش در سطح یک نمونه اثبات امکان‌پذیری محقق شده است. چارچوب توسعه‌یافته توانست مؤلفه‌های دوقلوی دیجیتال، عامل‌های تخصصی، مدل پیش‌بینی \gls{RUL} و عامل زبانی را در یک معماری واحد گردآوری کند و ارتباط میان این مؤلفه‌ها را در قالب فرایندهای پایش، تحلیل، تصمیم و کنترل نشان دهد. بنابراین، پاسخ کلی پژوهش به مسئله مطرح‌شده آن است که ترکیب این فناوری‌ها می‌تواند بستری برای ایجاد یک سامانه تصمیم‌یار تعاملی و قابل توسعه در زمینه نگهداری پیش‌بینانه فراهم کند.
در عین حال، نتایج نشان نمی‌دهند که چارچوب حاضر یک سامانه صنعتی نهایی یا یک مدل پیش‌بینی بهینه برای موتورهای توربوفن است. اعتبار نتایج پیش‌بینی به داده و شرایط آزمایش وابسته است و عملکرد عامل زبانی نیز در محدوده فرایندها و درخواست‌های تعریف‌شده ارزیابی شده است. ازاین‌رو، تعمیم مستقیم نتایج به محیط عملیاتی نیازمند اعتبارسنجی‌های گسترده‌تر، استفاده از داده‌های واقعی حسگر، بررسی شرایط متنوع عملیاتی و ارزیابی توسط متخصصان حوزه خواهد بود.
با وجود این، پژوهش حاضر نشان می‌دهد که ارزش چارچوب پیشنهادی صرفاً در دقت یک مدل پیش‌بینی منفرد خلاصه نمی‌شود، بلکه در نحوه قرار گرفتن این مدل در یک معماری عملیاتی و تعاملی است. در این معماری، پیش‌بینی می‌تواند به تشخیص وضعیت منجر شود، وضعیت می‌تواند مبنای تصمیم قرار گیرد، تصمیم می‌تواند به اقدام تبدیل شود و نتیجه اقدام می‌تواند دوباره در وضعیت دوقلوی دیجیتال بازتاب پیدا کند. این زنجیره، مبنایی برای توسعه نسل بعدی سامانه‌های نگهداری پیش‌بینانه فراهم می‌کند که در آن تحلیل داده، تصمیم‌گیری عامل‌محور و تعامل طبیعی کاربر در یک چارچوب مشترک قرار می‌گیرند.
در نهایت، می‌توان گفت که پژوهش حاضر یک گام عملی در جهت توسعه دوقلوهای دیجیتال چندعاملی تعاملی برای نگهداری پیش‌بینانه برداشته است. چارچوب ارائه‌شده از نظر معماری قابلیت توسعه دارد و می‌تواند در پژوهش‌های آینده با اتصال به داده‌های واقعی، افزایش سطح اعتبارسنجی، مدیریت صریح عدم‌قطعیت، تقویت سازوکارهای ایمنی و اعمال کنترل‌های انسانی در فرمان‌های حساس، به یک سامانه نزدیک‌تر به شرایط عملیاتی تبدیل شود.

% \section{نتیجه‌گیری}
% \label{sec:conclusion}
% هدف این پژوهش طراحی و پیاده‌سازی یک چارچوب دوقلوی دیجیتال چندعاملی برای پایش و پیش‌بینی عمر مفید باقیمانده موتورهای توربوفن و فراهم کردن امکان تعامل طبیعی با اطلاعات و وضعیت دوقلوی دیجیتال بود. برای دستیابی به این هدف، پژوهش از یک رویکرد سامانه‌محور استفاده کرد که در آن مدل پیش‌بینی، دوقلوی دیجیتال و عامل‌های کاربردی به صورت اجزای مستقل اما هماهنگ طراحی شدند. نتایج به‌دست‌آمده نشان می‌دهد که این رویکرد توانسته است یک نمونه اجرایی از ارتباط میان پیش‌بینی، پایش، تفسیر و کنترل ایجاد کند و در نتیجه، هدف اصلی پژوهش در سطح اثبات امکان‌پذیری معماری محقق شده است.
% در بخش پیش‌بینی، مدل \lr{V2} بر روی زیرمجموعه \lr{FD001} مجموعه داده \gls{CMAPSS} ارزیابی شد و مقدار \lr{MAE} برابر ۱۷٫۲۱ و مقدار $R^2$ برابر ۰٫۷۱ به دست آمد. همچنین نسبت کم‌برآورد از ۴۷ درصد در مدل پایه به ۸۷ درصد در مدل نهایی افزایش یافت. این نتایج نشان می‌دهد که مدل انتخاب‌شده توانایی فراهم کردن یک سرویس پیش‌بینی قابل استفاده در معماری را دارد و در عین حال، استفاده از تابع زیان نامتقارن توانسته است رفتار محافظه‌کارانه‌تری در برابر بیش‌برآورد \gls{RUL} ایجاد کند. با این حال، از آنجا که آزمایش تنها بر روی \lr{FD001} انجام شده است، این نتیجه نباید به عنوان اثبات برتری عمومی مدل یا تعمیم آن به تمامی شرایط عملیاتی تفسیر شود.
% در سطح معماری، نتایج نشان داد که تفکیک وظایف میان عامل‌ها امکان ایجاد یک زنجیره منسجم میان داده، پیش‌بینی، تفسیر و کنترل را فراهم کرده است. عامل پیش‌بینی وظیفه تخمین \gls{RUL} را انجام می‌دهد، عامل دوقلوی دیجیتال وضعیت موجودیت را نگهداری و ارائه می‌کند، عامل پایش وضعیت سامانه را مشاهده می‌کند و عامل زبان طبیعی اطلاعات حاصل را برای کاربر تفسیر کرده و در چارچوب مسیرهای تعریف‌شده، فرمان‌های مجاز را به اجزای مربوطه منتقل می‌کند. این جداسازی باعث می‌شود مسئولیت محاسبه عددی با مدل تخصصی باقی بماند و مدل زبان بزرگ به جای تولید مستقل مقدار \gls{RUL}، نقش واسط تعاملی و هماهنگ‌کننده را ایفا کند.
% یکی از نتایج مهم پژوهش، اجرای حلقه بسته میان پیش‌بینی، تصمیم، کنترل، به‌روزرسانی وضعیت و مشاهده مجدد است. این فرایند نشان می‌دهد که خروجی تحلیلی می‌تواند در یک معماری چندعاملی به یک اقدام عملیاتی متصل شود و نتیجه اقدام نیز مجدداً در وضعیت دوقلوی دیجیتال منعکس گردد. بنابراین، چارچوب پیشنهادی تنها یک سامانه برای نمایش یا پیش‌بینی وضعیت نیست، بلکه نمونه‌ای از یک سامانه تصمیم‌یار است که در آن اطلاعات حاصل از مدل یادگیری می‌تواند در یک فرایند عملیاتی قرار گیرد.
% با وجود این نتایج، باید میان اثبات امکان‌پذیری معماری و آمادگی برای استفاده صنعتی تمایز قائل شد. استفاده از داده‌های \gls{CMAPSS}، محدود بودن ارزیابی کمی به \lr{FD001}، شبیه‌سازی جریان داده به صورت شبه‌بی‌درنگ، وابستگی عامل زبان به سرویس بیرونی و نبود ارزیابی تخصصی مستقل برای همه توضیحات تولیدشده، عواملی هستند که دامنه نتیجه‌گیری را محدود می‌کنند. از این رو، چارچوب ارائه‌شده در وضعیت فعلی باید به عنوان یک نمونه آزمایشگاهی و اثبات امکان‌پذیری معماری در نظر گرفته شود، نه یک سامانه نهایی برای استقرار مستقیم در محیط عملیاتی.
% در مجموع، پاسخ پژوهش به مسئله مطرح‌شده را می‌توان چنین بیان کرد که ترکیب یک مدل پیش‌بینی تخصصی با دوقلوی دیجیتال، عامل‌های کاربردی و عامل تعامل زبانی می‌تواند یک مسیر عملی برای تبدیل پیش‌بینی عمر مفید از یک خروجی عددی مستقل به بخشی از یک سامانه یکپارچه پایش و تصمیم‌گیری فراهم کند. مزیت اصلی این رویکرد در ماژولار بودن معماری و امکان همکاری میان عامل‌های تخصصی است؛ به گونه‌ای که تغییر مدل پیش‌بینی، منبع داده یا روش تعامل لزوماً به بازطراحی کل سامانه منجر نمی‌شود. بنابراین، دستاورد اصلی پژوهش حاضر را می‌توان ارائه و ارزیابی یک معماری قابل توسعه دانست که در آن پیش‌بینی، بازنمایی دیجیتال، تعامل طبیعی و کنترل در یک چارچوب واحد به یکدیگر متصل شده‌اند.
% بر اساس این نتیجه، ادامه مسیر پژوهش باید بیش از آنکه صرفاً بر کاهش خطای یک مدل منفرد متمرکز باشد، بر افزایش قابلیت اعتماد، تعمیم‌پذیری و آمادگی عملیاتی کل سامانه تمرکز کند. انتقال ارزیابی از داده‌های شبیه‌سازی‌شده به داده‌های واقعی، بررسی شرایط عملیاتی متنوع‌تر، افزودن مدیریت پایدار وضعیت، کاهش وابستگی به سرویس‌های خارجی و ایجاد سازوکارهای کنترل دسترسی و تأیید انسانی، از جمله اقداماتی هستند که می‌توانند فاصله میان نمونه آزمایشگاهی حاضر و یک سامانه صنعتی را کاهش دهند. پیشنهادهای مشخص در این زمینه در بخش بعدی ارائه خواهد شد.


\section{پیشنهادهایی برای کارهای آینده}
\label{sec:future-work-final}

مسیرهای آینده پژوهش را می‌توان در چهار دسته اصلی تقسیم کرد.

در دسته نخست، باید دامنه داده گسترش یابد. آزمایش روی \lr{FD002} تا \lr{FD004} می‌تواند میزان مقاومت مدل و معماری را در برابر شرایط عملیاتی مختلف مشخص کند. همچنین استفاده از داده‌های واقعی حسگرهای موتور می‌تواند فاصله میان نمونه آزمایشگاهی و کاربرد صنعتی را کاهش دهد.

در دسته دوم، مدل پیش‌بینی می‌تواند از یک مدل نقطه‌ای به سمت مدل‌های احتمالاتی حرکت کند. ترکیب \lr{LSTM} با سایر مدل‌های احتمالاتی و استفاده از مدل‌های دقیق‌تر با خطای کم‌تر می‌تواند در دقت خروجی مدل کمک کند و در نتیجه خروجی زبان طبیعی را نیز واضح‌تر کند.

در دسته سوم، عامل زبانی می‌تواند با حافظه بلندمدت، پایگاه دانش فنی و دانش‌نگار حوزه هوافضا ترکیب شود. در این صورت، پاسخ عامل تنها بر اساس مقدار فعلی \gls{RUL} نخواهد بود، بلکه تاریخچه موتور، تعمیرات قبلی، شرایط عملیاتی و سوابق خرابی را نیز در نظر خواهد گرفت.

در دسته چهارم، معماری می‌تواند به حوزه‌های دیگر منتقل شود. برای مثال، در یک کارخانه تولیدی، عامل دوقلو می‌تواند نماینده یک دستگاه \lr{CNC} باشد؛ عامل پیش‌بینی وضعیت خرابی دستگاه را تخمین بزند؛ عامل پایش وضعیت خط تولید را بررسی کند و عامل زبانی امکان پرس‌وجو و اجرای فرمان را فراهم آورد. به همین ترتیب، معماری می‌تواند برای سایر صنایع مورد نیاز نیز توسعه یابد.

برای نزدیک شدن به استقرار صنعتی نیز پیشنهاد می‌شود ارتباط سرویس‌ها از مدل نظرسنجی دوره‌ای به سازوکارهای رویدادمحور منتقل شود و وضعیت دوقلو در یک زیرساخت پایدار ذخیره گردد. همچنین لازم است سیاست‌های کنترل دسترسی، ثبت رویداد، تأیید انسانی فرمان‌های حساس و سازوکارهای ایمنی در لایه عامل زبانی اضافه شوند.




% \section{جمع‌بندی نهایی}
% \label{sec:final-summary}

% هدف اصلی این پژوهش طراحی و پیاده‌سازی یک چارچوب دوقلوهای دیجیتال چندعاملی برای پایش و پیش‌بینی عمر مفید موتورهای توربوفن و بررسی امکان تعامل طبیعی با این دوقلوها بود. در این راستا، یک معماری چندلایه طراحی و سپس با استفاده از مجموعه داده \gls{CMAPSS} در قالب یک نمونه آزمایشگاهی پیاده‌سازی شد.

% نتایج نشان داد که مدل LSTM مورد استفاده می‌تواند به عنوان یک مؤلفه پیش‌بینی قابل استفاده در معماری عمل کند. مدل نهایی با استفاده از سازوکار توجه و \gls{PinballLoss} توانست \lr{MAE} برابر ۱۷٫۲۱ و $R^2$ برابر ۰٫۷۱ را روی مجموعه آزمون به دست آورد. همچنین نسبت کم‌برآورد از ۴۷ درصد در مدل پایه به ۸۷ درصد افزایش یافت. اهمیت نتیجه دوم در این است که مدل به صورت هدفمند به سمت پیش‌بینی محافظه‌کارانه حرکت کرده است؛ رویکردی که با ماهیت ایمنی‌محور نگهداری موتورها سازگارتر است. بااین‌حال، این نتایج به معنای ارائه بهترین مدل برای CMAPSS نیست و مدل پیش‌بینی در این پژوهش صرفاً به عنوان یک مؤلفه تخصصی و قابل جایگزینی در معماری در نظر گرفته شده است.

% در سطح سامانه‌ای، عامل زبانی توانست اطلاعات عددی تولیدشده توسط عامل پیش‌بینی را به زبان طبیعی تفسیر کند، وضعیت موتور را توضیح دهد و در فرایند‌های تعریف‌شده، فرمان‌های کنترلی را به دوقلوی دیجیتال ارسال کند. بررسی ۳۰ درخواست، نرخ موفقیت ۱۰۰ درصدی را در فرایند‌های تعریف‌شده نشان داد. همچنین حلقه بسته میان پیش‌بینی، تصمیم، کنترل و بازخورد در فرایند‌های عملیاتی اجرا شد.

% از دیدگاه معماری، مهم‌ترین دستاورد پژوهش در کنار هم قرار دادن قابلیت‌هایی است که معمولاً به صورت جداگانه بررسی می‌شوند: دوقلوی دیجیتال، سیستم چندعاملی، پیش‌بینی \gls{RUL}، پایش شبه‌بی‌درنگ و تعامل مبتنی بر مدل‌های بزرگ زبانی. این ترکیب امکان تبدیل یک مدل پیش‌بینی منفرد به یک سامانه تصمیم‌یار تعاملی را فراهم می‌کند. در چنین ساختاری، مدل یادگیری عمیق صرفاً تولیدکننده یک عدد نیست، بلکه خروجی آن وارد زنجیره‌ای می‌شود که می‌تواند به تشخیص وضعیت، ارائه توضیح، تصمیم‌گیری و در نهایت اقدام منجر شود.

% از منظر کاربرد صنعتی، این پژوهش یک نمونه اثبات امکان‌پذیری ارائه می‌کند و نه یک محصول آماده بهره‌برداری. برای رسیدن به سطح صنعتی، اتصال به داده‌های واقعی حسگر، پایگاه داده پایدار، مدیریت عدم‌قطعیت، اعتبارسنجی متخصصان، کنترل دسترسی و سازوکار تأیید انسانی فرمان‌های حساس ضروری است. با وجود این محدودیت‌ها، ساختار معماری به گونه‌ای طراحی شده است که این اجزا بتوانند در مراحل بعدی بدون تغییر بنیادی در منطق کلی سامانه اضافه شوند.

% در نهایت، می‌توان نتیجه گرفت که ارزش اصلی چارچوب پیشنهادی در \textbf{ترکیب پیش‌بینی، دوقلوی دیجیتال، عامل‌های تخصصی و تعامل زبانی در یک حلقه عملیاتی واحد} است. این ویژگی، معماری را از یک سامانه صرفاً پیش‌بینی‌کننده به سمت یک چارچوب دوقلوی دیجیتال چندعاملی قابل توسعه هدایت می‌کند؛ چارچوبی که می‌تواند در آینده با جایگزینی منبع داده، مدل پیش‌بینی و دانش حوزه، برای سایر تجهیزات و صنایع نیز مورد استفاده قرار گیرد.




